\begin{thebibliography}{99}
%\bibitem{ahoull} A.V. Aho and J.D. Ullman, Optimization of straight line programs, SIAM J. Comput. 1 (1972), 1--19.

\bibitem{alegur} N. Alechina and Yu. Gurevich, Syntax vs.\ semantics on finite structures, in \emph{Structures in Logic and Computer Science. A Selection of Essays in Honor of A. Ehrenfeucht}, Springer, 1997, pp. 14--33.

\bibitem{alm} J. Almeida, Finite Semigroups and Universal Algebra, World Scientific 1994.
%\bibitem{bgklp} L. Babai, A.J. Goodman, W.M. Kantor, E.M. Luks, P.P. P\'alfy, Short presentations for finite groups, \emph{J. Algebra}, {194} (1997), 79--112.

%\bibitem{bak} K. Baker, Finite equational bases for finite algebras in a congruence distributive equational class, Adv. Math. 24 (1977), 207--243.

\bibitem{BMW} K.A. Baker, G.F. McNulty and J. Wang,  An extension of Willard’s Finite Basis Theorem: Congruence meet-semidistributive varieties of finite critical depth, Algebra Universalis 52 (2005), 289--302.

%\bibitem{bakpix} K. A. Baker and A.F. Pixley,
%Polynomial interpolation and the Chinese remainder theorem for algebraic systems,
%Math. Z. 143 (1975), 165--174.

%\bibitem{hei} M. auf der Heide, A comparison of two variations of a pebble game on graphs, Theoretical Computer Science, 13 (1981), 315--322.

%\bibitem{babsze}  L. Babai and E. Szemer\'edi, On the complexity of matrix group problems I, in  Foundations of Computer Science, 1984. 25th Annual Symposium on Foundations of Computer Science. IEEE, 1984.

\bibitem{bel} V.P. Belkin, Quasi-identities of finite rings and lattices, Algebra Logika 17 (1978) 247--259 (in Russian).

\bibitem{ber} C. Bergman, Universal Algebra: Fundamentals and Selected Topics, Chapman Hall/CRC Press, 2011.

\bibitem{berslu00} C. Bergman and G. Slutzki, Complexity of some problems concerning varieties and quasi-varieties of algebras, SIAM J. Comput. 30 (2000), 359--382.


\bibitem{berslu00b}  C. Bergman and G. Slutzki, Computational complexity of some problems involving congruences on algebras. 15th Annual IEEE Symposium on Logic in Computer Science (Santa Barbara, CA, 2000), 168--174, IEEE Comput. Soc. Press, Los Alamitos, CA, 2000.

\bibitem{berslu02} C. Bergman and G. Slutzki, Computational complexity of some problems involving congruences on algebras, Theoret. Comp. Sci. 270 (2002),  591--608.

\bibitem{bez} G. Bezhanishvili, Locally finite varieties, Algebra Universalis 46 (2001), 531--548.

\bibitem{bir} G. Birkhoff, On the structure of abstract algebras, Math. Proc. Cambridge Phil. Soc.  31 (1935), 433--454.

%\bibitem{BKMS} A. Bulatov, M. Kozik, P. Mayr and M. Steindl, The subpower membership problem for semigroups, Internat. J. Algebra Comput. 26 (2016), 1435--1451.

\bibitem{BDJN} J. Bul\'{\i}n, D. Delic, M. Jackson and T Niven, A finer reduction of constraint problems to digraphs, Log. Methods Comp. Sci. 11 (2015), 1--33.

\bibitem{bursan} S. Burris and H.P. Sankappanavar, {A Course in Universal
Algebra}, Graduate Texts in Mathematics {78}, Springer Verlag, 1980.

\bibitem{CDFJ} D.M. Clark, B.A. Davey, R.S. Freese, M. Jackson, Standard topological algebras: Syntactic and principal congruences and profiniteness, Algebra Universalis 52 (2004) 343--376.

\bibitem{CDJP} D.M. Clark, B.A. Davey, M.G. Jackson and J.G. Pitkethly, The axiomatizability of topological prevarieties, {Adv. Math.} {218} (2008), 1604--1653.

\bibitem{clipre} Clifford A.H. and G.B. Preston, The algebraic theory
of semigroups, Vol. I, AMS Mathematical Surveys, (1961).

%\bibitem{cooset} S. Cook and R. Sethi, Storage requirements for deterministic polynomial time recognizable languages, J. Comput. System Sci. 13 (1976), 25--37.


\bibitem{cze} G. Cz\'edli, Cyclic congruences of slim semimodular lattices and non-finite axiomatizability of some finite structures, Arch. Math. (Brno) 58 (2022), 15--33.


\bibitem{DDK} J. Demel, M. Demlov\'a and V. Koubek, Fast algorithms constructing minimal subalgebras, congruences and ideals in a finite algebra, Theoret. Comp. Sci. 36 (1985), 203--216.

\bibitem{eil} S. Eilenberg, Automata languages and machines, Pure and Applied Mathematics, Volume B (Academic Press, New York, 1976).

\bibitem{eilsch} S. Eilenberg, M.-P. Sch\"utzenberger, On pseudovarieties,  Adv. Math. 19 (1976), 413--418.

\bibitem{ALV} R.S. Freese, R.N. McKenzie, G.F. McNulty and W.F. Taylor, Algebras, Lattices and Varieties, Volume II, Mathematical Surveys and Monographs 268, AMS 2022.

\bibitem{gor} V.A. Gorbunov, Algebraic Theory of Quasivarieties, Consultants Bureau, New York, 1998.

\bibitem{gur2} Y. Gurevich, Toward logic tailored for computational complexity, in Computation and Proof Theory, Lecture Notes in Mathematics 1104, 1984, pp. 175--216.

\bibitem{gussap} S.V. Gusev and O.B. Sapir, Classification of limit varieties of $\mathscr{J}$-trivial monoids, Comm. Algebra 50 (2002), 3007--3027.

\bibitem{ham} L. Ham, Relativised Homomorphism Preservation at the Finite Level, Studia Logica 105 (2017), 761--786.

%\bibitem{hewpat} C.E. Hewitt and M.S. Paterson, Comparative schematology, in ``Proceedings, Project MAC Conference on Concurrent Systems and Parallel Computation'', Woods Hole, Mass., 1970, pp. 119--127.

\bibitem{higjac} P.M. Higgins and M. Jackson, Algebras defined by equations, J. Algebra 555 (2020), 131--151.

\bibitem{higjac23} P.M. Higgins and M. Jackson, Equationally defined classes of semigroups, Semigroup Forum 107 (2023), 459--477.

\bibitem{hobmck} D. Hobby and R. McKenzie, The Structure of Finite Algebras, Contemporary Mathematics 76, American Mathematical Society, Providence, RI, 1988.

\bibitem{hod} W. Hodges, Model Theory,  Cambridge University Press, 1992.


\bibitem{imm} N. Immerman, Descriptive Complexity, Graduate Texts in Computer Science, Springer 1999.

%\bibitem{jac:INFB} M. Jackson, Small inherently nonfinitely based finite semigroups, Semigroup Forum 64 (2002), 297--324.

 \bibitem{jac:flat} M. Jackson, Flat algebras and the translation of universal Horn logic to equational logic, {J. Symbolic Logic} {73} (2008), 90--128.
 
 \bibitem{jac:SAT} M. Jackson, Flexible constraint satisfaction and a problem in semigroup theory, Internat. J. Algebra Comput. 35 (2025), 771--821.
 
 \bibitem{JKN} M. Jackson, T. Kowalski and T. Niven, Complexity and polymorphisms for digraph constraint problems under some basic constructions, Internat. J. Algebra Comput. 26 (2016), 1395--1433.
 
 \bibitem{jacmck} M. Jackson and R.N. McKenzie, Interpreting graph colorability in finite semigroups, Internat. J. Algebra Comput. 16 (2006), 119--140.
 
 \bibitem{jacmcn} M. Jackson and G.F. McNulty, The equational complexity of Lyndon’s algebra, Algebra Universalis 65 (2011), 243--262.
 
 \bibitem{JRZ} M. Jackson, M. Ren and X. Zhao, Nonfinitely based ai-semirings with finitely based semigroup reducts, J. Algebra 611 (2022), 211--245.

\bibitem{jactro10} M. Jackson and B. Trotta, The division relation: congruence conditions and axiomatisability, {Comm. Algebra} {38} (2010), 534--566.

\bibitem{jactro13} M. Jackson and B. Trotta, Constraint satisfaction,
irredundant axiomatisability
and continuous colouring, {Studia Logica} {101} (2013), 65--94.

\bibitem{jacvol} M. Jackson and M.V. Volkov, The Algebra of Adjacency Patterns: Rees Matrix Semigroups with Reversion, in A. Blass, N. Dershowitz, and W. Reisig (Eds.): Gurevich Festschrift, LNCS 6300, pp. 414--443, 2010.

%\bibitem{kalsch}  B. Kalyanasundaram and G. Schnitger, On the power of white pebbles, Combinatorica 11 (1991) 157--171.

\bibitem{kei} H.J. Keisler, Reduced products and Horn classes, Trans. Amer. Math. Soc. 117 (1965) 307--328.

\bibitem{kel} A.Kelarev, Graph Algebras and Automata, CRC Press, Boca Raton, 2003.


\bibitem{khasap} O.G. Kharlampovich and M.V. Sapir, Algorithmic problems in varieties, Internat. J. Algebra Comput. 5 (1995), 379--602.

\bibitem{KKP} O. Kl\'{\i}ma, M. Kunc and L. Pol\'{a}k, Deciding piecewise $k$-testability, manuscript (2014).

\bibitem{koz07} M. Kozik, Computationally and algebraically complex finite algebra membership problems, {Internat. J. Algebra Comput.} {17} (2007),  1635--1666.

\bibitem{koz09} M. Kozik, A 2EXPTIME complete varietal membership problem, {SIAM J. Comput.} {38} (2009),  2443--2467.

\bibitem{dag} A. Krokhin and S. \v{Z}vin\'y (eds), The Constraint Satisfaction Problem: Complexity and Approximability, Dagstuhl Follow Ups 7 (2017).

%\bibitem{lawwil}  J. Lawrence and R. Willard, On finitely based groups and nonfinitely based quasivarieties, J. Algebra 203 (1998), 1--11.

\bibitem{leezha} E.W.H. Lee and W.T. Zhang, Finite basis problem for semigroups of order six, LMS J. Comput.  Math. 18 (2015), 1--129.

\bibitem{lib} L. Libkin, Elements of Finite Model Theory. Texts in Theoretical Computer Science. An EATCS Series. Springer-Verlag, Berlin, 2004.

\bibitem{lyn} R.C. Lyndon, Identities in finite algebras, Proc. Amer. Math. Soc. 5 (1954), 8--9.

\bibitem{mck:lat} R. McKenzie, Equational bases for lattice theories,
Math. Scand. 27 (1970), 24--38.
\bibitem{mck1} R. McKenzie: The residual bounds of finite algebras. Internat. J. Alg. Comput. 6 (1996), 1--28. 

\bibitem{mck2} R. McKenzie: The residual bound of a finite algebra is not computable. Internat. J. Alg. Comput. 6 (1996), 29--48.

\bibitem{mck3} R. McKenzie: Tarski's finite basis problem is undecidable. Internat. J. Alg. Comput. 6 (1996), 49--104. 

%\bibitem{pal} P.P. P\'alfy, Unary polynomials in algebras. I., Algebra Universalis 18 (1984), 262--273.

\bibitem{par} R. Park, Equational classes of non-associative ordered systems, Ph.D. dissertation, UCLA, 1976.

\bibitem{pie} R. Pierce, Homomorphisms of semigroups, Ann. Math. 59 (1954), 287--291.

\bibitem{rei} O. Reingold, Undirected st-connectivity in log-space. In Harold N. Gabow and Ronald Fagin, editors, STOC, pages 376--385. ACM, 2005.

\bibitem{rei2} O. Reingold, Undirected connectivity in log-space, J. ACM 55(4) (2008), 1--24.

\bibitem{ros} B. Rossman, Homomorphism preservation theorems, {J. ACM} 55 (2008), no. 3, Art. 15, 53 pp.

\bibitem{rosen} E. Rosen, Some aspects of finite model theory and finite structures, Bull. Symb. Logic 8 (2002), 380--403.

%\bibitem{set} R. Sethi, Complete register allocation problems, SIAM J. Comput. 4 (1975), 226--248.

\bibitem{sap} M.V. Sapir, Sur la propri\'{e}t\'{e} de base finie pour les pseudovari\'{e}t\'{e}s de semigroupes finis. (French) [On the finite basis property for pseudovarieties of finite semigroups] C. R. Acad. Sci. Paris S\'{e}r. I Math. 306 (1988), 795--797.

\bibitem{shawil} S. Shaheen and R. Willard, Algebras from finite group actions and a question of Eilenberg and Sch\"utzenberger, J. Algebra (2023), doi: https://doi.org/10.1016/j.jalgebra.2023.05.014

\bibitem{sch66} B.M. Schein, Homomorphisms and subdirect decompositions of semigroups, Pacific J. Math. 17 (1966), 529--547.

\bibitem{sze} Z. Sz\'ekely, Computational complexity of the finite algebra membership problem for varieties, Internat. J. Algebra Comput. 12 (2002), 811--823.

\bibitem{tai} W.W. Tait, A counterexample to a conjecture of Scott and Suppes, J. Symbolic Logic 24 (1959), 15--16.

\bibitem{tum} V.I. Tumanov, On finite lattices having no independent
bases of quasi-Identities, Math. Notes 36 (1984), 625--634.

\bibitem{VGK} M.V. Volkov, S.V. Goldberg, and S.I. Kublanovsky, J. Math. Sci. 177 (2011), 847--859.

\bibitem{wil0} R. Willard, On McKenzie's method, Periodica Mathematica Hungarica 32 (1996), 149--165.

\bibitem{wil}
R. Willard, Tarski's finite basis problem
via ${\bf A}(\mathcal{T})$, {Trans.\ Amer.\ Math.\ Soc.} {349} (1997), 2755--2774.

\bibitem{wil2} R. Willard, A finite basis theorem for residually finite, congruence meet-semidistributive varieties, J. Symbolic Logic 65 (2000), 187--200.


\bibitem{wilFBP} R. Willard, The finite basis problem, Contributions to General Algebra, vol. 15, 199--206, Heyn, Klagenfurt, 2004.



\end{thebibliography}
